The electroweak link starts from the same angular operator that produced the vacuum constitutive relations, but uses another functional.\par

To obtain \(\mu\), \(\varepsilon\), \(Z\), and \(c\) it was necessary to preserve the full spectrum. It was necessary to distinguish the two sheets.\par

The determinant does the opposite:\par

\[
D_A=\det T=q_+q_-.
\]

By multiplying the eigenvalues, the determinant integrates the two orientations into a common invariant angular magnitude. It publishes the even part and reserves the leaf identity in the oriented memory.\par

By eso \(D_A\) is even under the inversion\par

\[
C^*\longmapsto-C^*.
\]

The sector \(W/Z\) has precisely that structure. The routes of \(W\) and \(Z\) are adjacent in the principal angular coordinate and share the same oriented coordinate. Upon forming the quotient, the contribution of \(C^*\) cancels.\par

Remains\par

\[
\left(\frac{m_W}{m_Z}\right)^2
=
D_A,
\]

and therefore\par

\[
\sin^2\theta_W
=
1-D_A.
\]

The identity reveals that the Weinberg angle is the complementary defect of the angular determinant.\par

The determinant measures the common part that survives when the two sheets are collapsed into their product. \(1-D_A\) measures the complementary part required for mixing.\par

It may be stated as follows: the vacuum preserves the entire response; the electroweak sector uses the even compression of that response.\par

The same angular antecedent produces two distinct functions:\par

\begin{itemize}[leftmargin=*,itemsep=0.35em]
\item the complete functional calculus generates the constitutive relations;
\item the determinant generates the mixing invariant.
\end{itemize}

This explains why the vacuum and the Weinberg angle are two distinct functional realizations of the same angular antecedent.\par

From that point onward, the couplings \(g\) and \(g'\) are, at the stated level, the two projections of the same rationalized electromagnetic coupling onto the sine and cosine directions of the mixing:\par

\[
g^2
=
\frac{4\pi\alpha_{\rm HMT}}{1-D_A},
\qquad
g'^2
=
\frac{4\pi\alpha_{\rm HMT}}{D_A}.
\]

The tree custodial identity\par

\[
\rho_{\rm EW}^{(0)}=1
\]

states that the same angular partition organizes mass and coupling coherently.\par

The Fermi constant appears later as the effective low-energy response of the charged channel:\par

\[
G_F^{(0)}
=
\frac{\pi\alpha_{\rm HMT}}
{\sqrt2(1-D_A)m_W^2}.
\]

The Fermi constant appears as a composition of the fine structure, the angular defect, and the carrier scale \(W\).\par

The Higgs sector introduces a decisive difference. Its route has an oriented coordinate distinct from the route of \(W\), so \(C^*\) persists. The scalar self-coupling publishes the orientation reserved by the determinant.\par

The gauge sector is even: it uses \(D_A\).\par
The scalar sector is oriented: it retains \(C^*\).\par

This division is physically suggestive. Gauge fields organize the common part of the transformation; the scalar sector preserves the memory of the broken orientation.\par

Radiative corrections belong to a later level and must be constructed prospectively from the internal loops, states, and scales, with \(G_F\) as output and terminal evaluation. The causal skeleton is already present: vacuum, mixing, couplings, and Fermi response follow from different readings of the same enriched angular state.\par

